% 2026-11-16 站主問：「本文中的『協調』？」「協調的價值與 Alexander 與 Schauer 的定紛止爭有何不同」。
% Law 757--785 頁由第二路 agent 抽出原句，758、762、784 頁主對話回文字層核過；Alexander 與 Schauer 頁碼取自 ../W09_預習/Teile/20_AS.tex（2026-11-09 已核）。
\chapter{協調與定紛止爭}
\label{ch:60_coordination}

\begin{qaquestion}
Law\ix{persons}{law@Law, David S.} 第三部以協調\ix{terms}{01@\zhlidxcat{賽局理論}!協調賽局}說明法院的權力；W09 Alexander\ix{persons}{alexander@Alexander, Larry} 與 Schauer\ix{persons}{schauer@Schauer, Frederick} 以定紛止爭\ix{terms}{04@\zhlidxcat{定紛止爭論}!定紛止爭}說明其他部門應服從最高法院的解釋。
\begin{enumerate}
\tightlist
\item 本文所說的協調指什麼？協調的對象有哪幾種？
\item 協調與定紛止爭有何相同、有何不同？
\end{enumerate}
\end{qaquestion}

兩說都以一個人人遵從的答案為價值，也都不以該答案正確為前提。不同在於：Alexander 與 Schauer 給的是服從的規範理由，對象是其他官員，目的在穩定；Law 說明的是人民與政府事實上何時服從，對象擴及全體人民，而同一種協調也可以用來推翻政府。

\section{本文中的協調}
\label{sec:60_coordination-1}

\begin{enumerate}
\tightlist
\item
  定義。「In many situations, people have an interest in cooperating with one another, yet they cannot cooperate effectively unless they can coordinate their actions」〔許多情形下，人們有彼此合作的利益，卻非協調其行動不能有效合作〕；這類兼有共同利益與衝突利益的賽局，Schelling 稱為 mixed-motive games\ix{terms}{01@\zhlidxcat{賽局理論}!混合動機賽局}〔混合動機賽局〕，一切議價情境皆屬之。訴訟也是如此：兩造偏好的結果不同，卻都從在法庭上解決爭議而得益。Law 以划船隊的舵手為例：每個划手都知道其他划手聽到了舵手的口令，也預期他們照口令划；法庭在訴訟當事人之間造成同樣的信念。\footcite[758]{law2009}
\item
  協調的機制是改變人們對他人行為的預期。757 頁區分法院塑造的三種信念：規範信念、事實信念、預測信念；協調靠的是第三種，「By generating predictive beliefs about the behavior of others, courts can induce rational actors to comply with their decisions」〔藉由造成對他人行為的預測信念，法院得以使理性行為者遵從其裁判〕。\footcite[757]{law2009}預期要成為共同知識\ix{terms}{01@\zhlidxcat{賽局理論}!共同知識}，判決須是顯著而唯一的焦點\ix{terms}{01@\zhlidxcat{賽局理論}!焦點}\pg{770--775 頁}，且須夠具體\pg{777--778 頁}。\footcite[770--778]{law2009}
\item
  協調的對象有三種。
  \begin{enumerate}
  \tightlist
  \item 訴訟兩造之間。邊界爭端的兩國都寧可讓步也不願開戰，判決選定一個均衡\pg{759--760 頁}；歐洲法院在 Cassis de Dijon\ix{cases}{03@歐洲共同體!01@歐洲法院!1979Cassis de Dijon（1979）@Cassis de Dijon（1979）} 案以相互承認原則提供成員國協調的焦點，Law 說「Coordination, not coercion, is the more apt description of what courts do in mixed games of conflict and cooperation」〔在兼有衝突與合作的賽局中，協調比強制更貼切地描述法院所為〕\pg{761--762 頁}。
  \item 全體公民對判決的服從。George v. Albert 假設案中，只要人人相信他人都知道判決、人人相信他人服從法院、人人相信抵抗一個人人服從的判決既昂貴又徒勞，法院不論判決有無說服力或受不受歡迎，都可以把全體公民協調到它所選的結果上\pg{763--764 頁}。
  \item 人民對抗政府。推翻暴政須眾人同時行動，人人不動也是均衡；法院判政府違憲，即是起事所需的公開信號\pg{765--770 頁}。政權的存續取決於人們對他人信念的信念，「The power to change people's beliefs about how others will behave is, in short, the power to topple governments」〔改變人們對他人行為之信念的權力，就是推翻政府的權力〕\pg{768 頁}。
  \end{enumerate}
  \footcite[759--770]{law2009}
\item
  協調與權力。法院的權力取決於判決是否被服從，與判決受不受歡迎無關\pg{780 頁}；不受歡迎而仍被服從的判決越多，法院的協調能力越強\pg{780--782 頁}。\footcite[780--782]{law2009}
\end{enumerate}

\section{與定紛止爭的比較}
\label{sec:60_coordination-2}

\begin{longtable}{@{}L{0.16\linewidth}L{0.38\linewidth}L{0.38\linewidth}@{}}
\caption{Law 的協調與 Alexander 與 Schauer 的定紛止爭。}\\
\toprule
項目 & Law：協調 & Alexander 與 Schauer：定紛止爭 \\
\midrule
\endfirsthead
\toprule
項目 & Law：協調 & Alexander 與 Schauer：定紛止爭 \\
\midrule
\endhead
所問 & 法院無兵無財，人民與政府何以服從（757、800--801 頁） & 其他官員是否應將最高法院的解釋視為與憲法本身同等權威（455 頁） \\
服從的理由 & 各方預期他人服從，服從即為最佳回應（763--764 頁） & 權威提供與內容無關之理由，長期而言以定紛為準比逐案自行判斷犯錯較少（457--458 頁註 11） \\
被協調或受拘束者 & 訴訟兩造、全體公民、對抗政府的人民 & 其他部門官員，司法與非司法皆同（455 頁） \\
判決有誤時 & 人人預期他人遵守的錯誤判決，同樣是均衡 & 解釋可能有誤，不減其權威（457--458 頁） \\
效果的方向 & 可維持秩序，也可協調人民推翻政府（768 頁） & 穩定、可預測、使決定得以作成、行動得以協調（467 頁） \\
法院以外 & 憲法、領袖、重大事件亦可作為協調裝置，法院的優勢在顯著與唯一（772--775 頁） & 法院是否為最佳的定紛者，另於第 II 部分 B 節討論（467--478 頁） \\
\bottomrule
\end{longtable}

\begin{enumerate}
\tightlist
\item
  兩說的共同點在於答案本身的價值。Alexander 與 Schauer 說，法律可以提供「neither final nor conclusive, but nevertheless authoritative」〔既非終局、也非結論性，但仍有權威〕的定紛，對一階歧見提供二階的解決，使決定得以作成、行動得以協調。\footcite[467]{alexander2000}Law 的焦點同樣不問判決是否正確，只問它是否被人人預期為他人所遵守。兩說都用到「協調」一詞，Alexander 與 Schauer 的協調是定紛的效果，Law 的協調是法院權力的來源。
\item
  兩說回答的問題不同。Alexander 與 Schauer 457 頁說，最高法院在 Cooper\ix{cases}{02@美國!01@聯邦最高法院!1958Cooper v. Aaron（1958）@Cooper v. Aaron（1958）} 與 City of Boerne 宣告自身的至上地位，與問題無關，因為這種宣告預設了其他部門與公眾是否應接受的結論。\footcite[457]{alexander2000}Law 784 頁註 198 引 Cooper 宣告聯邦司法於憲法解釋居於至上地位的那段，並引兩人 1997 年的前文；他敘述 Cooper 作成於艾森豪派遣第 101 空降師進駐小岩城之後，法院在服從已有保障時才宣告其至上地位，「With the passage of time, the illusion of power has become the reality」〔時間一久，權力的幻象成了現實〕\pg{785 頁}。\footcite[784--785 \& n.198]{law2009}同一則宣告，在 Alexander 與 Schauer 是有待證成的規範主張，在 Law 是因被服從而成真的預期。
\item
  兩說的受拘束者不同。Alexander 與 Schauer 處理的是其他部門官員應否服從最高法院；Law 的協調對象主要是人民，政府官員服從法院，是因為預期人民會依判決行動\pg{794 頁}。\footcite[794]{law2009}定紛止爭在部門之間運作，協調在人民與政府之間運作。
\item
  兩說對穩定的態度不同。Alexander 與 Schauer 以穩定與可預測為定紛的價值；Law 的協調可以維持服從均衡，也可以使人民從默從的均衡移到起事的均衡，判決由此成為推翻政府的信號\pg{768--769 頁}。定紛止爭的論證無法說明法院判政府違憲時人民何以起而反抗，Law 的協調論則把協調人民對抗政府列為法院的主要功能。
\end{enumerate}
